\begin{figure}[p]
\centering
\includegraphics[width=\textwidth]{supplementary/figure_measured_alignment_power.pdf}
\caption{\textbf{Conditional sensitivity of the measured ancestry-response test.}
\textbf{A}, Detection of positive alignment as the fraction of reliable tuning variance carried by a shared-path covariance increases. Every simulated dataset preserves the thirteen scans, seven targets, observed partners, stimulus identities and repeats. Green uses measured repeat reliabilities; gray removes measurement noise.
\textbf{B}, The same detection probabilities versus the mean simulated seven-target partial-rank effect. This is a model-conditional effect axis, not a universal detectable correlation. The dashed line in \textbf{A,B} marks 80\%; bands are Wilson 95\% Monte Carlo intervals from 4,000 paired complete datasets. Detection requires an exact two-sided seven-target signed-rank $P\leq0.05$ and a positive mean effect.
\textbf{C}, Mean split-half Spearman reliability in 1,000 independent calibration datasets versus the observed value, for 125 partner records. The two negative estimates are assigned zero reliable variance (orange squares); the dotted segment shows that clipping. The identity line is not fitted. All partner-pair correlations come from joint response matrices; pairs are not inferential replicates. See Supplementary Note~\ref{note:measured_alignment_power} for the injected covariance and calibration assumptions.}
\label{fig:supp_measured_alignment_power}
\end{figure}
